% AA-27 | Iteracja XIV; szczegółowy audyt historyczny III, 20 XI 2026.
\chapter[Rozszerzenia Artina--Schreiera]{Rozszerzenia Artina--Schreiera w charakterystyce dodatniej}
\label{aa:ch:artin-schreier}

\section{Wielomian addytywny i jego pierwiastki}

W rozdziale~\ref{aa:ch:kummer} do opisu pewnych rozszerzeń cyklicznych
wykorzystano równania potęgowe. Jeżeli jednak $F$ ma charakterystykę $p$,
to pochodna wielomianu $X^p-a$ jest zerowa. Wielomian ten nie może zatem
pełnić tej samej roli w teorii rozszerzeń rozdzielczych stopnia $p$.
Zamiast niego rozważymy $X^p-X-a$: dodatkowy składnik $-X$ zapewnia
niezerową pochodną, a tożsamość $(x+y)^p=x^p+y^p$ pozwala opisać
wszystkie pierwiastki i automorfizmy.

\begin{definicja}\label{aa:def:as-map}
Niech $p$ będzie liczbą pierwszą, a $F$ ciałem charakterystyki $p$.
Przez $\mathbb F_p\subseteq F$ oznaczamy ciało pierwsze. Odwzorowanie
\(\wp_F:(F,+)\longrightarrow(F,+),\quad \wp_F(x)=x^p-x\),
nazywamy \emph{odwzorowaniem Artina--Schreiera}. Jego obraz oznaczamy
przez $\wp(F)=\{x^p-x:x\in F\}$. Wszystkie rozszerzenia algebraiczne
$F$ rozpatrujemy we wspólnym, ustalonym domknięciu algebraicznym
$\overline F$.
\end{definicja}
\begin{proof}[Uzasadnienie własności odwzorowania]
Dla dowolnych $x,y\in F$ dwumian Newtona daje
\((x+y)^p=x^p+y^p\): dla $1\le j\le p-1$ współczynnik
$\binom pj$ jest podzielny przez $p$, a więc jest zerem w $F$.
W konsekwencji $\wp_F(x+y)=\wp_F(x)+\wp_F(y)$, co dowodzi
addytywności. Jeżeli $\wp_F(x)=0$, to $x$ jest pierwiastkiem
$X^p-X$. Każdy $c\in\mathbb F_p$ spełnia $c^p=c$:
dla $c=0$ jest to natychmiastowe, a dla $c\ne0$ wynika z faktu,
że grupa $\mathbb F_p^\times$ ma rząd $p-1$. Wielomian stopnia
$p$ nie ma więcej niż $p$ pierwiastków w ciele. Zatem
$\ker\wp_F=\mathbb F_p$. W szczególności $\wp(F)$ jest podgrupą
addytywnej grupy $F$, a iloraz $F/\wp(F)$ ma sens.
\end{proof}

\begin{twierdzenie}[Konstrukcja rozszerzenia Artina--Schreiera]\label{aa:tw:as-forward}
Niech $p$ będzie liczbą pierwszą, $F$ ciałem charakterystyki $p$,
$a\in F$, zaś $\alpha\in\overline F$ pierwiastkiem
$f_a(X)=X^p-X-a$. Jeżeli $a\notin\wp(F)$, to $f_a$ jest
nierozkładalny nad $F$, rozszerzenie $E=F(\alpha)/F$ jest Galois
stopnia $p$, a przyporządkowanie
\(c\longmapsto\sigma_c,\quad \sigma_c(\alpha)=\alpha+c\)
wyznacza izomorfizm grupy addytywnej $(\mathbb F_p,+)$ na
$\operatorname{Gal}(E/F)$.
\end{twierdzenie}
\begin{proof}
\emph{Krok 1: pierwiastki i rozdzielczość.} Ponieważ
$f_a'(X)=pX^{p-1}-1=-1$, wielomian $f_a$ nie ma pierwiastków
wielokrotnych. Dla każdego $c\in\mathbb F_p$ zachodzi
\((\alpha+c)^p-(\alpha+c)=\alpha^p-\alpha+c^p-c=a\).
Elementy $\alpha+c$ są parami różne, gdyż dodawanie $\alpha$
jest iniekcją. Jest ich dokładnie $p$, czyli tyle, ile wynosi
stopień $f_a$. Wszystkie należą do $E$, więc $E$ jest ciałem
rozkładu rozdzielczego wielomianu $f_a$ nad $F$. Stąd $E/F$
jest skończonym rozszerzeniem Galois.

\emph{Krok 2: zanurzenie grupy Galois.} Niech
$G=\operatorname{Gal}(E/F)$ i $\sigma\in G$.
Ponieważ $\sigma$ ustala współczynniki $f_a$,
element $\sigma(\alpha)$ jest pewnym pierwiastkiem $f_a$.
Istnieje zatem dokładnie jeden $c_\sigma\in\mathbb F_p$,
dla którego $\sigma(\alpha)=\alpha+c_\sigma$.
Odwzorowanie $G\to(\mathbb F_p,+)$ dane przez
$\sigma\mapsto c_\sigma$ jest iniektywne: jeśli
$c_\sigma=c_\tau$, to $\sigma$ i $\tau$ zgadzają się na
$\alpha$ i na $F$, a zatem na całym $F(\alpha)$. Jest także
homomorfizmem, ponieważ dla $\sigma,\tau\in G$
\[
(\sigma\circ\tau)(\alpha)
=\sigma(\alpha+c_\tau)
=\alpha+c_\sigma+c_\tau,
\]
gdzie użyliśmy faktu $c_\tau\in\mathbb F_p\subseteq F$.
W szczególności $|G|$ dzieli $p$, gdyż $G$ jest izomorficzna
z podgrupą grupy addytywnej rzędu $p$.

\emph{Krok 3: eliminacja przypadku trywialnego.}
Jeżeli $|G|=1$, to $[E:F]=|G|=1$, ponieważ $E/F$ jest Galois.
Wtedy $\alpha\in F$, skąd $a=\alpha^p-\alpha\in\wp(F)$,
wbrew założeniu. Zatem $|G|=p$ i $[E:F]=p$.
Wielomian minimalny $\alpha$ nad $F$ dzieli $f_a$ i ma
stopień $[F(\alpha):F]=p$, wobec czego jest równy $f_a$.
Zanurzenie $G\to(\mathbb F_p,+)$ ma teraz obraz rzędu $p$,
a więc jest izomorfizmem. Odwrotne przyporządkowanie daje
wskazane automorfizmy $\sigma_c$.
\end{proof}

\begin{twierdzenie}[Odwrotne twierdzenie Artina--Schreiera]\label{aa:tw:as-converse}
Niech $p$ będzie liczbą pierwszą, $F$ ciałem charakterystyki $p$
i $E/F$ rozszerzeniem Galois, którego grupa jest cykliczna rzędu
$p$. Wówczas istnieją $\alpha\in E$ i $a\in F\setminus\wp(F)$
takie, że $E=F(\alpha)$ oraz $\alpha^p-\alpha=a$.
\end{twierdzenie}
\begin{proof}
Niech $\sigma$ będzie generatorem $G=\operatorname{Gal}(E/F)$.
Ponieważ $E/F$ jest rozdzielcze, jego forma śladu
$(x,y)\mapsto\operatorname{Tr}_{E/F}(xy)$ jest niezdegenerowana
(zob. rozdział~\ref{aa:ch:trace-norm}). W szczególności sam
funkcjonał $x\mapsto\operatorname{Tr}_{E/F}(x)$ nie jest
zerowy: w przeciwnym razie forma śladu byłaby identycznie
zerowa. Jest to niezerowy funkcjonał do jednowymiarowej
$F$-przestrzeni $F$, więc istnieje $u\in E$ o śladzie $1$
(po pomnożeniu dowolnego elementu o niezerowym śladzie przez
odwrotność tego śladu).

Definiujemy
\[
\alpha=-\sum_{j=0}^{p-1}j\,\sigma^j(u),
\]
przy czym $j$ oznacza obraz liczby całkowitej $j$ w
$\mathbb F_p\subseteq F$. Sprawdzimy zmianę $\alpha$ pod
$\sigma$, zamiast przyjmować ją bez uzasadnienia. Po zmianie
indeksu $k=j+1$ oraz użyciu $\sigma^p=\mathrm{id}$ mamy
\[
\sigma(\alpha)=-\sum_{k=1}^{p-1}(k-1)\sigma^k(u)
-(p-1)u,
\qquad
\alpha=-\sum_{k=1}^{p-1}k\sigma^k(u).
\]
Odejmowanie daje
\[
\sigma(\alpha)-\alpha
=\sum_{k=1}^{p-1}\sigma^k(u)-(p-1)u
=\sum_{k=0}^{p-1}\sigma^k(u)
=\operatorname{Tr}_{E/F}(u)=1,
\]
ponieważ $-(p-1)=1$ w charakterystyce $p$.
Następnie
\[
\sigma(\alpha^p-\alpha)
=(\alpha+1)^p-(\alpha+1)=\alpha^p-\alpha.
\]
Zatem $a:=\alpha^p-\alpha$ należy do ciała elementów
stałych $E^G=F$. Ponieważ $\sigma(\alpha)\ne\alpha$,
element $\alpha$ nie należy do $F$. Z drugiej strony
$[E:F]=p$ i $F\subsetneq F(\alpha)\subseteq E$, więc z
multiplikatywności stopni wynika $[F(\alpha):F]=p$ i
$E=F(\alpha)$. Gdyby $a=b^p-b$ dla pewnego $b\in F$,
to $(\alpha-b)^p-(\alpha-b)=0$. Własność jądra $\wp_F$
dowiedziona wyżej (w domknięciu algebraicznym również)
pokazuje, że $\alpha-b\in\mathbb F_p\subseteq F$; otrzymalibyśmy
$\alpha\in F$, sprzeczność. Stąd $a\notin\wp(F)$.
\end{proof}

\section{Przykłady, kontrprzykłady i klasyfikacja}

\begin{przyklad}[Rozszerzenie ciała funkcji wymiernych]\label{aa:prz:as-rational}
Niech $F=\mathbb F_p(t)$, gdzie $t$ jest zmienną
transcendentną nad $\mathbb F_p$. Uzasadnimy, że
$t\notin\wp(F)$. Każdy niezerowy $b\in F$ zapisuje się
jako $P(t)/Q(t)$, gdzie $P,Q\in\mathbb F_p[t]$ są niezerowe.
Niech $m=\deg P-\deg Q$, czyli rząd bieguna
$b$ w punkcie nieskończoności, jeśli $m>0$.
Gdy $m>0$, funkcja $b^p$ ma stopień w nieskończoności $pm$,
ściśle większy od $m$; dlatego $b^p-b$ również ma tam
biegun rzędu $pm$. Gdy $m\le0$, zarówno $b^p$, jak i $b$
nie mają bieguna w nieskończoności, więc również ich różnica
nie ma takiego bieguna. Tymczasem $t$ ma biegun rzędu $1$.
Równość $b^p-b=t$ jest niemożliwa, ponieważ $pm\ne1$
dla żadnej liczby całkowitej $m>0$. Dla $b=0$ jest to
tym bardziej jasne. Twierdzenie~\ref{aa:tw:as-forward}
daje więc cykliczne rozszerzenie Galois
$\mathbb F_p(t,\alpha)/\mathbb F_p(t)$ stopnia $p$, gdzie
$\alpha^p-\alpha=t$. Jego automorfizmy mają postać
$\alpha\mapsto\alpha+c$ dla $c\in\mathbb F_p$.
\end{przyklad}

\begin{przyklad}[Dlaczego warunek na klasę $a$ jest konieczny?]\label{aa:prz:as-trivial}
Niech $b\in F$ i $a=b^p-b$. Wówczas dla każdego
$c\in\mathbb F_p$ mamy $(b+c)^p-(b+c)=a$, skąd
\[
X^p-X-a=\prod_{c\in\mathbb F_p}(X-b-c).
\]
Wszystkie czynniki są różne, bo $b+c=b+d$ oznacza $c=d$.
Wielomian pozostaje więc rozdzielczy, lecz jego ciało rozkładu
to samo $F$. Na przykład nad $F=\mathbb F_2(t)$ dla
$b=t$ otrzymujemy $a=t^2+t$ i rozkład
$X^2+X+t^2+t=(X+t)(X+t+1)$.
\end{przyklad}

\begin{stwierdzenie}[Zależność od klasy addytywnej]\label{aa:prop:as-classes}
Niech $F$ będzie ciałem charakterystyki $p$ i $a,b\in F$.
Jeżeli istnieją $c\in\mathbb F_p^\times$ oraz $d\in F$
takie, że $b=ca+\wp_F(d)$, to ciała rozkładu wielomianów
$X^p-X-a$ i $X^p-X-b$ w ustalonym $\overline F$ są równe.
\end{stwierdzenie}
\begin{proof}
Wybierzmy $\alpha\in\overline F$ o
$\alpha^p-\alpha=a$ i połóżmy $\beta=c\alpha+d$.
Ponieważ $c^p=c$ dla $c\in\mathbb F_p$, otrzymujemy
$\beta^p-\beta=c(\alpha^p-\alpha)+d^p-d=b$.
Wszystkimi pierwiastkami pierwszego wielomianu są
$\alpha+\mathbb F_p$, a drugiego $\beta+\mathbb F_p$.
Ich ciała rozkładu to odpowiednio $F(\alpha)$ i
$F(\beta)$. Z równości $\beta=c\alpha+d$ wnioskujemy
$F(\beta)\subseteq F(\alpha)$. Ponieważ $c\ne0$,
mamy też $\alpha=c^{-1}(\beta-d)\in F(\beta)$;
zatem zachodzi równość.
\end{proof}

\begin{twierdzenie}[Klasyfikacja rozszerzeń cyklicznych stopnia $p$]
\label{aa:tw:as-classification}
Niech $F$ będzie ciałem charakterystyki $p$, a $a,b$ dwoma
elementami $F\setminus\wp(F)$. Niech $E_a,E_b\subset\overline F$
będą ciałami rozkładu odpowiednio $X^p-X-a$ i $X^p-X-b$.
Wówczas $E_a=E_b$ wtedy i tylko wtedy, gdy istnieją
$c\in\mathbb F_p^\times$ i $d\in F$ takie, że
$b=ca+\wp_F(d)$. Innymi słowy, jednowymiarowe
podprzestrzenie przestrzeni wektorowej $F/\wp(F)$ nad
$\mathbb F_p$ odpowiadają bijektywnie cyklicznym
rozszerzeniom Galois $F$ stopnia $p$ wewnątrz $\overline F$.
\end{twierdzenie}
\begin{proof}
Wystarczalność dowiedziono w
stwierdzeniu~\ref{aa:prop:as-classes}. Dla konieczności
załóżmy $E_a=E_b=:E$. Wybierzmy pierwiastki
$\alpha^p-\alpha=a$ i $\beta^p-\beta=b$ w $E$.
Na mocy twierdzenia~\ref{aa:tw:as-forward} istnieje generator
$\sigma$ grupy $\operatorname{Gal}(E/F)$ spełniający
$\sigma(\alpha)=\alpha+1$. Ponieważ $\sigma$ jest
$F$-automorfizmem, element $\sigma(\beta)$ także jest
pierwiastkiem $X^p-X-b$, a więc
$\sigma(\beta)=\beta+c$ dla pewnego $c\in\mathbb F_p$.
Gdyby $c=0$, to $\beta\in E^{\langle\sigma\rangle}=F$,
co dawałoby $b\in\wp(F)$. Zatem $c\ne0$.
Teraz $\sigma(\beta-c\alpha)=\beta-c\alpha$, więc
$d:=\beta-c\alpha\in F$. Z tożsamości
$\beta=c\alpha+d$ wynika
$b=\beta^p-\beta=c(\alpha^p-\alpha)+d^p-d
=ca+\wp_F(d)$.
Na koniec $\wp(F)$ jest podprzestrzenią $F$ nad
$\mathbb F_p$ (addytywność oraz
$\wp_F(cx)=c\wp_F(x)$ dla $c\in\mathbb F_p$).
Zatem warunek $b=ca+\wp_F(d)$ jest dokładnie równością
prostych $\mathbb F_p[a]=\mathbb F_p[b]$ w ilorazie,
gdzie $[a]$ oznacza klasę $a$. Konstrukcja z pierwszego
twierdzenia i jego odwrotność zapewniają istnienie oraz
pełność tej klasyfikacji.
\end{proof}

\begin{przyklad}[Charakterystyka dwa i ciało czteroelementowe]
\label{aa:prz:as-four-elements}
Dla $F=\mathbb F_2$ mamy $\wp_F(0)=\wp_F(1)=0$,
więc $1\notin\wp(F)$. Wielomian $X^2-X-1=X^2+X+1$
nie ma pierwiastka w $\mathbb F_2$ i wyznacza
$E=\mathbb F_2[\alpha]$ o czterech elementach.
W $E$ zachodzi $\alpha^2=\alpha+1$, dlatego
$\sigma(\alpha)=\alpha^2=\alpha+1$ jest nietrywialnym
$F$-automorfizmem. Ponieważ
$\sigma^2(\alpha)=(\alpha+1)^2=\alpha^2+1=\alpha$,
jest to generator grupy cyklicznej rzędu $2$.
W tym przypadku $\alpha^2-\alpha=1$ i wszystkie formuły
Artina--Schreiera można bezpośrednio sprawdzić.
\end{przyklad}

\begin{uwaga}\label{aa:uw:as-vs-kummer}
Warunek $a\notin\wp(F)$ oznacza, że klasa $[a]$ jest
niezerowa w $F/\wp(F)$. W charakterystyce $p$ wielomian
$X^p-a$ ma pochodną zero i, gdy ma pierwiastek w domknięciu
algebraicznym, jest potęgą $(X-\gamma)^p$.
Natomiast $X^p-X-a$ ma pochodną $-1$, niezależnie od $a$.
To wyjaśnia, dlaczego w opisie \emph{rozdzielczych}
rozszerzeń cyklicznych stopnia $p$ stosujemy konstrukcję
Artina--Schreiera, a nie równanie potęgowe z teorii Kummera.
\end{uwaga}
